\section{Related Work}
\paragraph{Hallucination detection.} Internal-state probes
\cite{azaria2023internal} and contrast-consistent search
\cite{burns2023discovering} both rely on supervised access to
internal features. Our approach uses an entirely structural
descriptor of the attention graph, with no learned parameters.

\paragraph{Topological data analysis (TDA).} Persistent homology
\cite{edelsbrunner2008persistent} provides a multi-scale topological
summary robust to noise. Applications to neural networks include
analyses of loss landscapes \cite{horoi2022exploring} and feature
spaces \cite{naitzat2020topology}. To our knowledge no published
work measures topology of the attention graph as a hallucination
signal.

\paragraph{Attention graphs.} Treating attention as a graph has a
long history~\cite{vig2019analyzing}, but typically for visualisation;
quantitative graph invariants are uncommon.
